% Part: incompleteness
% Chapter: incompleteness-provability
% Section: lob-thm

\documentclass[../../../include/open-logic-section]{subfiles}

\begin{document}

\olfileid{inc}{inp}{tar}

\olsection{سچائی دی ناقابلِ تعریفیت}

\emph{قابلِ تعریف ہون} دا تصور حساب دی زبان لئی رسمی معنیات
ہون اُتے منحصر اے۔ اسیں حساب دی زبان دے فارمولیاں تے جملیاں
دا سیٹ بیان کیتا اے۔ «مرادی تعبیر» وچ ایہہ جملے قدرتی
اعداد بارے دعوے کردے نیں،تے ایسا دعویٰ سچ یا جھوٹ ہو سکدا
اے۔ $\Struct{N}$ اوہ !!{structure} ہووے جیہدا دائرہ
$\Nat$ اے تے حساب دی زبان دیاں علامتاں دی معیاری تعبیر
ورتدا اے۔ فیر $\Sat{N}{!A}$ دا مطلب اے: «$!A$ معیاری
تعبیر وچ سچ اے»۔

\begin{defn}
قدرتی اعداد دا کوئی تعلق $R(x_1,\dots,x_k)$، $\Struct{N}$
وچ اُدوں تے صرف اُدوں \emph{قابلِ تعریف} اے جدوں حساب دی زبان وچ
ایسا فارمولا $!A(x_1,\dots,x_k)$ ہووے کہ ہر
$n_1,\dots,n_k$ لئی $R(n_1,\dots,n_k)$ اُس ویلے تے
صرف اوس ویلے درست ہووے جدوں $\Sat{N}{!A(\num n_1,\dots,\num
  n_k)}$۔
\end{defn}

دوجے لفظاں وچ، کوئی تعلق $\Struct{N}$ وچ اُدوں تے صرف اُدوں
قابلِ تعریف اے جدوں اوہ نظریے $\Th{TA}$ وچ قابلِ نمائندگی
ہووے؛ایتھے $\Th{TA} =
\Setabs{!A}{\Sat{N}{!A}}$ حساب دے سچ جملیاں دا سیٹ
اے۔ (جے ایہہ گل فوراً واضح نہ ہووے،تاں تعریفاں ول
مڑ کے پڑتال کرو تے اپنے آپ نوں اس دی وجہ سمجھاؤ۔)

\begin{lem}
ہر قابلِ حساب تعلق~$\Struct{N}$ وچ قابلِ تعریف اے۔
\end{lem}

\begin{proof}
آسانی نال پڑتال ہو جاندی اے کہ $\Th{Q}$ وچ کسے تعلق
دی نمائندگی کرن والا فارمولا، $\Struct{N}$ وچ اوہی تعلق
متعین کردا اے۔
\end{proof}

ہن پُچھ سکدے آں کہ کی الٹ وی درست اے؟ یعنی
$\Struct{N}$ وچ قابلِ تعریف ہر تعلق قابلِ حساب اے؟
جواب نہیں۔ مثال لئی:

\begin{lem}
حساب رُکّن دا تعلق $\Struct{N}$ وچ قابلِ تعریف اے۔
\end{lem}

\begin{proof}
یاد کرو کلینی دے نارمل روپ دا قضیہ کہندا اے کہ ہر جزوی
قابلِ حساب فنکشن~$f$ دا کوئی اشاریہ~$e$ اے جیہڑے نال
$f(x) =
U(\umin{s}{T(e,x,s)})$ ہر $x \in \Nat$ لئی؛ایتھے
$U$ تے $T$ بنیادی بازگشتی
تے ایس لئی کُل نیں۔ سو $f(x)$ اُدوں تے صرف اُدوں متعین اے، یعنی حساب اُدوں ہی رُکدا
اے، جدوں کوئی ایہو جیہا~$s$ ہووے جیہڑے لئی
$T(e,x,s)$ درست ہووے۔

ہن $H$ حساب رُکّن دا تعلق ہووے،یعنی
\[
H = \Setabs{\tuple{e,x}}{\lexists[s][T(e, x, s)]}.
\]
$!D_T$، $T$ دی $\Struct{N}$ وچ تعریف کردا ہووے۔ فیر
\[
H = \Setabs{\tuple{e,x}}{\Sat{N}{\lexists[s][!D_T(\num e, \num x, s)]}},
\]
ایس لئی $\lexists[s][!D_T(z, x, s)]$،
~$H$ نوں $\Struct{N}$ وچ تعریف کردا اے۔
\end{proof}

\begin{prob}
دکھاؤ کہ $Q(n) \defiff n \in \Setabs{\Gn{!A}}{\Th{Q} \Proves !A}$
حساب وچ قابلِ تعریف اے۔
\end{prob}

خود $\Th{TA}$ دا کیہہ؟ کی اوہ حساب وچ قابلِ تعریف اے؟
یعنی کی سیٹ $\Setabs{\Gn{!A}}{\Sat{N}{!A}}$ حساب
وچ قابلِ تعریف اے؟ تارسکی دے قضیے دا جواب ناں وچ اے۔

\begin{thm}
\ollabel{thm:tarski}
حساب دے سچ !!{sentence}s دا سیٹ حساب وچ قابلِ تعریف
نہیں۔
\end{thm}

\begin{proof}
منّو $!D(x)$ اُس دی تعریف کردا اے؛یعنی
$\Sat{N}{!A}$ اُدوں تے صرف اُدوں جدوں
$\Sat{N}{!D(\gn{!A})}$۔ مقرر نقطے دے لمے نال کوئی
فارمولا $!A$ ملدا اے جیہڑے لئی
$\Th{Q} \Proves !A \liff \lnot !D(\gn{!A})$،
تے ایس توں $\Sat{N}{!A \liff \lnot !D(\gn{!A})}$۔
پر فیر $\Sat{N}{!A}$ اُدوں تے صرف اُدوں جدوں
$\Sat{N}{\lnot !D(\gn{!A})}$،جو ایس مفروضے دے
خلاف اے کہ $!D(y)$ حساب دے سچ بیانواں دا سیٹ تعریف
کردا اے۔
\end{proof}

تارسکی نے ایہہ تجزیہ سچائی دے ہور عمومی فلسفیانہ
تصور اُتے وی لاگو کیتا۔ کوئی وی زبان $L$ منّو؛تارسکی
دے مطابق $L$ لئی سچائی دا مناسب تصور ہر جملے $X$ واسطے
ایہہ شرط پوری کرے گا:
\begin{quote}
«$X$» اُدوں تے صرف اُدوں سچ اے جدوں $X$۔
\end{quote}
انگریزی وچ تارسکی دی بار بار نقل ہون والی مثال
“Snow is white” اے؛پنجابی وچ ایہہ شرط ایویں پڑھیے:
\begin{quote}
«برف چٹی اے» اُدوں تے صرف اُدوں سچ اے جدوں برف چٹی اے۔
\end{quote}
پر ہر ایسی زبان وچ جیہڑی قطری فنکشن دی نمائندگی لئی کافی
قوی ہووے،تے ہر لسانی محمول $T(x)$ لئی،اسیں ایسا جملہ
$X$ بنا سکدے آں جیہڑا ایہہ آکھے: «$X$ اُدوں تے صرف اُدوں جدوں
$T(\text{«$X$»})$ نہیں»۔ چونکہ اسیں نہیں چاہندے کہ
سچائی دا محمول کسے جملے نوں سچ تے جھوٹ دونوں قرار دے،
تارسکی نے نتیجہ کڈھیا کہ ایسی زبان دے سارے جملیاں لئی
سچائی دا محمول متعین کرن لئی کسے نہ کسے طرح اوس زبان
دیاں حداں توں باہر جانا پئے گا۔ دوجے لفظاں وچ،ایسی
زبان دا سچائی محمول خود اوسے زبان وچ قابلِ تعریف نہیں۔

\end{document}
